\documentclass[b5j,fleqn]{jarticle}
\usepackage{zennyakuUniv}
\begin{document}
\body{
\Title{newcrown3 p72,p73}
\para

\sent{What is "Fairness"?}
{「公平」とは何でしょうか。}

\para

\sent{Fairness is about treating people equally, and we all want life to be fair.}
{公平とは人々を平等に扱うことであり、私たちはみな、人生が公平であってほしいと願っています。}

\sent{How can we be fair to other people?}
{私たちはどうすれば、ほかの人に公平に接することができるでしょうか。}

\para

\sent{There are many ways to be fair.}
{公平であるための方法はたくさんあります。}

\sent{One way is to give everyone the same resources and opportunities.}
{一つの方法は、全員に同じ資源と機会を与えることです。}

\sent{For example, when you are distributing four cookies to two children, how many will you give each?}
{例えば、4枚のクッキーを2人の子どもに分けるとき、それぞれに何枚ずつ与えますか。}

\sent{The most common answer is "two" because each child can get the same number of cookies.}
{最も一般的な答えは「2枚」です。なぜなら、それぞれの子どもが同じ枚数のクッキーをもらえるからです。}

\sent{What will you do if there are five cookies?}
{クッキーが5枚あったら、どうしますか。}

\sent{If you give the extra cookie to one child, it will seem unfair.}
{余ったクッキーを一人の子どもに与えると、不公平に思えるでしょう。}

\sent{So you will probably break the cookie into two pieces and give one to each child.}
{そこでおそらく、クッキーを二つに割り、それぞれの子どもに一つずつ与えるでしょう。}

\para

\sent{Another way to be fair is to give each person the exact resources and opportunities they need.}
{公平であるためのもう一つの方法は、一人ひとりに、その人が必要とする資源と機会を必要なだけ与えることです。}

\sent{Here are two pictures that explain it.}
{それを説明する二つの絵を見てみましょう。}

\para

\sent{In the first picture, each person is given one box to watch a baseball game.}
{最初の絵では、野球の試合を見るために、一人ひとりに箱が一つずつ与えられています。}

\sent{Everyone gets the same number of boxes.}
{全員が同じ数の箱を受け取ります。}

\sent{However, the result is not the same.}
{しかし、結果は同じではありません。}

\sent{The child on the right cannot see the game with just one box.}
{右側の子どもは、箱が一つだけでは試合を見ることができません。}

\para

\sent{In the next picture, each person gets a different number of boxes.}
{次の絵では、一人ひとりが異なる数の箱を受け取ります。}

\sent{The child on the right gets two boxes, the child in the middle gets one box, and the man on the left gets no box.}
{右側の子どもは箱を二つ、真ん中の子どもは一つ受け取り、左側の男性は箱を受け取りません。}

\sent{They can all see the game.}
{これで全員が試合を見ることができます。}

\sent{When you think about fairness in this way, you recognize differences and consider each person's needs.}
{このように公平について考えるとき、私たちは違いを認識し、一人ひとりの必要性を考慮します。}

\para

\sent{Fairness can be achieved in many ways.}
{公平は、さまざまな方法で実現できます。}

\sent{Think about the best way to be fair in different situations.}
{さまざまな状況で公平であるための最善の方法を考えてみましょう。}
}
\end{document}
